\chapter{Runtime Profiles and Anytime Performance}
\label{cha:profiles}



\section{Motivation}

Section~\ref{subsec:bbo-perf} surveyed the three standard frameworks for measuring
the performance of black-box optimization algorithms: the expected runtime
$\mathbb{E}[\tau_{f^*}]$, fixed-budget analysis, and fixed-target analysis culminating
in the runtime profile framework from \cite{pinto2018practice}. Each
improves on its predecessor in descriptive richness, yet none provides a single scalar
that captures solution quality at every point in the run. Fixed-budget analysis, for
instance, compresses the full trajectory to one value per budget~\cite{jansen2014fixed}.
What the practitioner actually observes is the evolution of solution quality over the
\emph{entire} course of the run.
 
Pinto and Doerr make the case for this richer mode of reporting directly:
\begin{displayquote}
Our intention is rather to highlight to researchers working on theoretical aspects in
evolutionary computation that, beyond being possibly more relevant for practitioners,
explicitly reporting runtime profiles can give quite interesting insights into the
performance of \acrshort{ea}s.
{(\cite{pinto2018practice}, p.~19)}
\end{displayquote}

This chapter takes that framing as its starting point. Building on the notation established in Section~\ref{subsubsec:fixed-target}, this chapter formalizes the distributional runtime profile as a mathematical
object, and then proves that expected cumulative regret decomposes exactly as a
double sum over this profile. The result unifies the runtime-analysis and
regret-based perspectives into a single identity, providing the
\emph{combined anytime measure} \citet{pinto2018practice} call for.
 


% \section{Fixed-Budget Analysis}

% As introduced in Section~\ref{subsubsec:fixed-budget}, fixed-budget analysis
% characterizes $\mathbb{E}[f^{*}_T(\mathcal{A})]$ (the expected quality of the best
% solution found by algorithm $\mathcal{A}$ within a budget of $T$ evaluations) and
% establishes that maximizing this quantity is equivalent to minimizing expected simple
% regret (Section~\ref{subsubsec:sr}) at budget $T$~\cite{jansen2014fixed}. Fixed-budget
% analysis therefore furnishes a natural starting point: it converts the binary "found /
% not found" question of expected runtime into a quality-valued statement, but still
% compresses the entire trajectory into a single scalar per budget value. The runtime
% profile framework, introduced below, lifts this to a trajectory-level characterization.


\section{Variants}

\subsection{Expectation-Based Runtime Profiles}

Building on the first hitting time $\tau_v$ defined in
Section~\ref{subsubsec:fixed-target}, \citet{pinto2018practice} propose
reporting the full collection of expected first hitting times
$\bigl(\mathbb{E}[\tau_v]\bigr)_{v}$ across all relevant target levels, the
\textit{expectation-based runtime profile}. When canonical fitness levels exist, such
as for \textsc{OneMax} or \textsc{LeadingOnes}, these provide natural targets; for
other functions, intermediate values identified during analysis or by linear
interpolation of the fitness range may be used~\cite{pinto2018practice}. No new
algorithms or proof techniques are required: the profile is typically already implicit
in existing runtime proofs and calls only for a richer mode of presentation.
 
The practical value of this richer view is illustrated by the behavior of $(1+1)$~\acrshort{ea}
variants on \textsc{LeadingOnes}: reporting only the expected optimization time suggests
RLS is superior, yet the runtime profile reveals that
$(1+1)$~EA$_{>0}$ (\textit{App. para.}~\ref{app:par:resampea})
reaches all intermediate fitness levels $i \leq 7{,}980$ faster than RLS for problem
size $n = 10{,}000$, with RLS only pulling ahead in the final portion of the
run~\cite{pinto2018practice}. Such crossover behavior is directly relevant to the
design of adaptive parameter and operator selection schemes.

\subsection{Distributional Runtime Profiles}

The expectation-based runtime profile of~\cite{pinto2018practice} summarizes the
distribution of $\tau_v$ by its mean. A finer characterization, and the appropriate
quantity for connecting runtime profiles to regret-based measures, is the
\textit{distributional runtime profile}
\begin{equation}
    \mathcal{P}(v, t) \;:=\; \mathbb{P}(\tau_v > t),
\end{equation}
which records, for each target level $v$ and time step $t$, the probability that level
$v$ has not yet been reached by time $t$. The expectation-based profile is recovered by
the standard tail-sum formula for non-negative integer-valued random
variables~\cite[p.~16]{doerr2020theory}:
\begin{equation}
    \label{eqn:dist-runtime-rel}
    \mathbb{E}[\tau_v] \;=\; \sum_{t=0}^{\infty} \mathbb{P}(\tau_v > t).
\end{equation}
The distributional profile is the more primitive object from which expectation-based
statements follow, and it constitutes the formal link between runtime profiles and
the cumulative regret framework developed below.

\begin{remark}[Naming convention]
\label{rem:profile-naming}
The term \emph{runtime profile} appears in two distinct senses in the literature and in
this work. \cite{pinto2018practice} uses it to denote the
\emph{expectation-based} runtime profile $\bigl(\mathbb{E}[\tau_v]\bigr)_v$, which
summarizes per-level performance by expected hitting times. The present work additionally
employs the strictly finer \emph{distributional} runtime profile $\mathcal{P}(v,t) =
\mathbb{P}(\tau_v > t)$, from which the expectation-based profile is recovered via the
tail-sum formula (Equation~\ref{eqn:dist-runtime-rel}). To avoid ambiguity, both
notions are referred to explicitly whenever both are in scope. Hereafter, and in all
practical experimentation in this report, \emph{runtime profile} refers to the
distributional variant $\mathcal{P}(v,t)$ unless otherwise stated.
\end{remark}

\section{Connection to Cumulative Regret}
\label{subsec:regret-profiles}

One of the directions suggested by Pinto and Doerr in their \textit{Future Works}
section is restated for emphasis:
\begin{displayquote}
    We feel that there is a need for combined measures that are capable of describing
    the anytime behavior of an \acrshort{ea}. Regret-based measure as
    used in machine learning could be a key here, but we haven't been able so far to
    identify a fully satisfying measure.
    {(\cite{pinto2018practice}, p.~22)}
\end{displayquote}
As far as we know, establishing such a relationship remains an open question.
This section consolidates the following proposal: that the \textit{\acrfull{cr}}
(Section~\ref{subsubsec:cr}) and the distributional runtime profile are not independent
approaches to algorithm performance but are two representations of the same underlying
quantity. The identity below attempts to establish this relation.

The relationship can be made explicit through a formal setup. A randomized black-box
algorithm produces an adaptive sequence of query points $x_1, x_2, \ldots \in
\{0,1\}^n$. Recall from Section~\ref{subsubsec:cr} the \textit{best-so-far process}
\[
    B(t) \;:=\; \max_{s \leq t} f(x_s),
\]
which records the best objective value seen up to time $t$. Since $B(t)$ is
non-decreasing in $t$, the event $\{B(t) < v\}$ is equivalent to $\{\tau_v > t\}$:
the level $v$ has not yet been hit. The incumbent cumulative regret over a budget of
$T$ evaluations is therefore
\[
    \mathrm{CR}(T) \;:=\; \sum_{t=1}^{T} \bigl(f^* - B(t)\bigr).
\]

The proof of the main theorem relies on two elementary lemmas.

\begin{lemma}[Tail-sum representation for the gap]
\label{lem:tailsum}
For any integer $b$ with $0 \le b \le f^*$,
\[
    f^* - b \;=\; \sum_{v=1}^{f^*} \mathbf{1}[b < v].
\]
\end{lemma}

\begin{proof}
The right-hand side counts the integers $v \in \{1,\ldots,f^*\}$ that strictly exceed
$b$, namely the set $\{b+1,\ldots,f^*\}$, which has cardinality $f^* - b$.
\end{proof}

\begin{lemma}[Indicator equivalence]
\label{lem:indicator}
For all $v \in \{1, \ldots, f^*\}$ and $t \in \mathbb{N}$,
\[
    \mathbf{1}[B(t) < v] \;=\; \mathbf{1}[\tau_v > t].
\]
\end{lemma}

\begin{proof}
Since $B(t)$ is non-decreasing in $t$, the event $\{B(t) \ge v\}$ is equivalent to
$\{\tau_v \le t\}$. Taking complements gives $\{B(t) < v\} = \{\tau_v > t\}$, so the
two indicators agree on every sample path.
\end{proof}

The following identity connects \textit{expected} cumulative regret to the distributional
runtime profile.

\begin{theorem}[Tail-Sum Regret Representation]
\label{thm:regret-profile}
For any randomized black-box algorithm operating on $f$ with finite budget
$T \in \mathbb{N}$,
\[
    \mathbb{E}[\mathrm{CR}(T)]
    \;=\;
    \sum_{v=1}^{f^*} \sum_{t=1}^{T} \mathbb{P}(\tau_v > t).
\]
\end{theorem}

\begin{proof}
For each $t$, Lemma~\ref{lem:tailsum} applied to $b = B(t)$ gives
\[
    f^* - B(t) \;=\; \sum_{v=1}^{f^*} \mathbf{1}[B(t) < v].
\]
Summing over $t = 1, \ldots, T$ yields
\[
    \mathrm{CR}(T)
    \;=\; \sum_{t=1}^{T}\bigl(f^* - B(t)\bigr)
    \;=\; \sum_{t=1}^{T} \sum_{v=1}^{f^*} \mathbf{1}[B(t) < v].
\]
By Lemma~\ref{lem:indicator}, $\mathbf{1}[B(t) < v] = \mathbf{1}[\tau_v > t]$ on every
sample path, so
\[
    \mathrm{CR}(T) \;=\; \sum_{t=1}^{T} \sum_{v=1}^{f^*} \mathbf{1}[\tau_v > t].
\]
Taking expectations and applying linearity of expectation together with
$\mathbb{E}[\mathbf{1}[\tau_v > t]] = \mathbb{P}(\tau_v > t)$ gives
\[
    \mathbb{E}[\mathrm{CR}(T)]
    \;=\; \sum_{t=1}^{T} \sum_{v=1}^{f^*} \mathbb{P}(\tau_v > t).
\]
Since both sums are finite, their order may be exchanged, yielding the stated result.
\end{proof}

Theorem~\ref{thm:regret-profile} expresses expected cumulative regret as a double sum
of the distributional runtime profile: $\mathbb{E}[\mathrm{CR}(T)]$ is the total
expected time the algorithm spends below each fitness level $v$, summed over all levels.
An algorithm with a uniformly smaller runtime profile (one that reaches each fitness
level earlier in expectation) therefore incurs lower cumulative regret. The regret
measure thus simultaneously rewards early progress at every tier of the landscape,
formalizing the intuition that steady improvement throughout a run is preferable to
stagnation followed by a late breakthrough, even when both strategies share the same
expected optimization time.

Taking the infinite-horizon limit, and applying the tail-sum identity for
$\mathbb{E}[\tau_v]$ along with the monotone convergence theorem, yields the following
asymptotic connection.

\begin{corollary}[Asymptotic Regret and Expected Hitting Times]
\label{cor:asymptotic-regret}
Suppose that $\tau_v$ is almost surely finite with finite expectation for each
$v \in \{1, \ldots, f^*\}$. Then
\[
    \lim_{T \to \infty} \mathbb{E}[\mathrm{CR}(T)]
    \;=\;
    \sum_{v=1}^{f^*} \mathbb{E}[\tau_v] \;-\; f^*.
\]
\end{corollary}
 
\begin{proof}
\textbf{Step 1: Tail-sum identity for $\mathbb{E}[\tau_v]$.}
Since $\tau_v \in \mathbb{N}_{\geq 1}$ by the convention that $t = 1$ is the first
oracle call and initialization is excluded from the query budget, we have
$\mathbb{P}(\tau_v > 0) = 1$. The standard tail-sum formula for a positive
integer-valued random variable then gives
\[
    \mathbb{E}[\tau_v]
    \;=\; \sum_{t=0}^{\infty} \mathbb{P}(\tau_v > t)
    \;=\; 1 + \sum_{t=1}^{\infty} \mathbb{P}(\tau_v > t),
\]
hence $\sum_{t=1}^{\infty} \mathbb{P}(\tau_v > t) = \mathbb{E}[\tau_v] - 1$.
 
\textbf{Step 2: Passing to the limit in Theorem~\ref{thm:regret-profile}.}
For each fixed $v$, the partial sums $\sum_{t=1}^{T}\mathbb{P}(\tau_v > t)$ are
non-decreasing in $T$ and non-negative, so the monotone convergence theorem gives
\[
    \lim_{T\to\infty} \sum_{t=1}^{T} \mathbb{P}(\tau_v > t)
    \;=\; \sum_{t=1}^{\infty} \mathbb{P}(\tau_v > t)
    \;=\; \mathbb{E}[\tau_v] - 1.
\]
 
\textbf{Step 3: Summing over $v$.}
Applying Step 2 to each $v = 1,\ldots, f^*$ and invoking
Theorem~\ref{thm:regret-profile},
\[
    \lim_{T \to \infty} \mathbb{E}[\mathrm{CR}(T)]
    \;=\; \sum_{v=1}^{f^*}\bigl(\mathbb{E}[\tau_v] - 1\bigr)
    \;=\; \sum_{v=1}^{f^*} \mathbb{E}[\tau_v] \;-\; f^*. \qedhere
\]
\end{proof}
 
Corollary~\ref{cor:asymptotic-regret} provides a unifying perspective: the long-run
expected cumulative regret decomposes as the sum of the expected first hitting times
over all fitness levels, minus the constant $f^*$ arising from the $t = 0$ term in the
tail-sum identity. Algorithms that reach all intermediate fitness levels quickly are
therefore those that incur low long-run cumulative regret. This formalizes
the intuition advanced in~\cite{pinto2018practice} that an algorithm making steady
progress through the fitness landscape is preferable to one that stagnates at low
fitness values even if both ultimately find the optimum in the same expected time.
 
The identity also reveals that the regret-based measures studied in the online learning
literature~\cite{cesa2006prediction} and the runtime profile framework
of~\cite{pinto2018practice} are not independent approaches to measuring algorithm
performance but are two views of the same underlying quantity, connected through the
distributional runtime profile $\mathcal{P}(v, t)$.

\section{Practical Relevance}
\label{subsec:profiles-relevance}

As discussed in Section~\ref{subsubsec:perf-limitations}, the limitations of
classical performance measures stem from their treatment of optimization as a
sequence of threshold events rather than as a continuous quality trajectory.

Pinto and Doerr~\cite{pinto2018practice} note that runtime profiles are particularly
valuable for the design of non-static parameter and operator selection strategies in
\acrshort{ea}s (such as adaptive mutation rates, crossover probabilities, and selection
pressure) since such strategies must make decisions based on current solution quality
rather than on the expected optimization time alone. The identification of
\emph{crossover fitness levels} (where one algorithm overtakes another in the runtime
profile) provides actionable guidance for when to switch between strategies.

Furthermore, as established in Theorem~\ref{thm:regret-profile}, any collection of
upper bounds on $\mathbb{E}[\tau_v]$ for all levels $v$ simultaneously yields an upper
bound on expected cumulative regret, and conversely, lower bounds on regret propagate to
lower bounds on a weighted aggregate of first hitting times. Cumulative regret thus
provides a single scalar summary of the entire runtime profile, complementing the
per-level view that the profile itself offers. The two representations (regret as a
scalar, runtime profile as a function) are therefore best understood as complementary
instruments for the same measurement task, connected by the identity of
Theorem~\ref{thm:regret-profile}.